Squaring the norm shows that the six oriented differences among $a,b,c$, reduced modulo $7$, must satisfy
\[
1+\sum_{\substack{u,v\in\{a,b,c\}\\u\ne v}}\zeta^{u-v}=0.
\]
The minimal polynomial $1+x+\cdots+x^6$ of $\zeta$ implies that each nonzero residue must appear exactly once among those differences: the displayed polynomial has degree at most $6$, constant coefficient $1$, and its only possible multiple of the minimal polynomial is the polynomial itself.

Equivalently, the three cyclic gaps between the chosen residues are $1,2,4$ in some order. Indeed the unordered partitions of $7$ into three positive gaps are $(1,1,5),(1,2,4),(1,3,3),(2,2,3)$, and only $(1,2,4)$ produces distinct distance classes $1,2,3$. There are two cyclic orders of these three distinct gaps and seven translations, giving fourteen triples in $\mathbb Z_7$. Each residue belongs to $14\cdot\dfrac{3}{7}=6$ of them, so eight avoid $0$ and are eligible here.

Negating residues pairs each eligible triple of sum $s$ with one of sum $21-s$. Hence the total is $8(\dfrac{21}{2})=\boxed{84}$.
